\documentclass[12pt,a4paper]{ctexart}
\usepackage[margin=1in]{geometry}
\usepackage{booktabs}
\usepackage{array}
\usepackage{amsmath}
\usepackage{tabularx}
\usepackage{xurl}
\usepackage[colorlinks=true,linkcolor=blue,urlcolor=blue]{hyperref}
\newcolumntype{L}[1]{>{\raggedright\arraybackslash}p{#1}}

\title{Step5a 算法与实验准备（中文版，T32 主口径）}
\author{}
\date{\today}

\begin{document}
\maketitle

\section{文档目的}
本稿用于组会汇报，统一当前实验口径，并给出可直接写入论文的结果表述。
\begin{enumerate}
  \item 数据协议统一为 \textbf{table\_group + train:val=1:5 + 无测试集}；
  \item 主文消融统一在 \textbf{TCDR(32,64)} 上进行；
  \item TCDR(16,96) 仅作为补充对照，不作为主文主表。
\end{enumerate}

\section{数据协议}
\begin{itemize}
  \item 划分模式：\texttt{table\_group}
  \item 划分比例：\textbf{train:val = 1:5}
  \item 测试集：\textbf{不使用（test=0）}
\end{itemize}

统一数据集：
\begin{center}
\texttt{output/dataset\_table\_group\_train\_val\_no\_test.pkl}
\end{center}
目标域样本统计：train=1274，val=6270，test=0。

\section{模型概述（TCDR）}
共享主干先提取特征：
\begin{equation}
\mathbf{h}=f_{\theta}([\mathbf{x};\mathbf{z}]).
\end{equation}
共享线性头给出基预测：
\begin{equation}
\hat{y}_{base}=\mathbf{w}^{\top}\mathbf{h}+b.
\end{equation}
再引入拓扑条件残差专家：
\begin{equation}
\Delta_{src}=g_{\phi_{src}}([\mathbf{h};\mathbf{e}_t]),\quad
\Delta_{tgt}=g_{\phi_{tgt}}([\mathbf{h};\mathbf{e}_t]),
\end{equation}
最终预测：
\begin{equation}
\hat{y}=
\begin{cases}
\hat{y}_{base}+\Delta_{src}, & \text{source}\\
\hat{y}_{base}+\Delta_{tgt}, & \text{target}
\end{cases}
\end{equation}

\section{统一训练设置}
主实验固定参数（除被消融项）：
\begin{itemize}
  \item \texttt{--freeze\_hgat}
  \item \texttt{--num\_epoch 260}, \texttt{--batch\_size 2048}
  \item \texttt{--learning\_rate 2e-3}, \texttt{--enc\_lr\_scale 0.05}
  \item \texttt{--weight\_decay 5e-5}, \texttt{--mlp\_dropout 0.2}
  \item \texttt{--hgat\_l2\_norm}, \texttt{--z\_noise\_std 0.01}
  \item \texttt{--hgat\_par\_cap\_weight 0.6}
  \item \texttt{--src\_loss\_anneal\_start 80}, \texttt{--src\_loss\_anneal\_end 280}
  \item \texttt{--src\_loss\_final\_scale 0.8}
  \item Scheduler=\texttt{plateau}，factor=0.5，patience=3，threshold=3e-4，min\_lr=1e-6
  \item Early stop：patience=12，min\_delta=8e-4
  \item \texttt{--val\_eval\_interval 5}, \texttt{--skip\_test\_eval}
\end{itemize}

\subsection{评价指标与统计规则}
为避免不同 seed 最优轮次不一致导致的统计偏差，本稿统一采用日志中的 \texttt{[Best]} 行（由验证损失选择）进行聚合，统计列为：
\begin{itemize}
  \item \texttt{Train R\textsuperscript{2}@Best}；
  \item \texttt{Val R\textsuperscript{2}}；
  \item \texttt{Val Loss}（MSE）；
  \item \texttt{Best Epoch}。
\end{itemize}
其中
\begin{equation}
R^2 = 1-\frac{\sum_i (y_i-\hat{y}_i)^2}{\sum_i (y_i-\bar{y})^2},\qquad
\text{MSE} = \frac{1}{N}\sum_i (y_i-\hat{y}_i)^2 .
\end{equation}
汇总结果均报告为 5-seed 的 mean$\pm$std。

\section{核心结果（Core，多 seed）}
来源：\texttt{paper\_materials/tables/paper\_notest\_core\_summary.csv}

\begin{table}[htbp]
\centering
\footnotesize
\setlength{\tabcolsep}{4pt}
\begin{tabular}{L{3.3cm}c c c c}
\toprule
方法 & NSeed & Train $R^2$ (mean$\pm$std) & Val $R^2$ (mean$\pm$std) & Val Loss (mean$\pm$std) \\
\midrule
Baseline(shared linear) & 5 & 0.9024$\pm$0.0086 & 0.8726$\pm$0.0067 & 0.121635$\pm$0.006428 \\
TCDR(16,96) & 5 & 0.9586$\pm$0.0085 & 0.9065$\pm$0.0168 & 0.089249$\pm$0.016086 \\
TCDR(32,64) & 5 & \textbf{0.9582$\pm$0.0216} & \textbf{0.9234$\pm$0.0175} & \textbf{0.073105$\pm$0.016755} \\
\bottomrule
\end{tabular}
\caption{Core 架构对比（1:5，无测试集）。}
\label{tab:core_cn}
\end{table}

\subsection{Core 结果解读}
\begin{enumerate}
  \item 相比 Baseline(shared linear)，TCDR(32,64) 的 Val $R^2$ 绝对提升 \textbf{+0.0508}（0.8726 $\rightarrow$ 0.9234），Val Loss 相对下降约 \textbf{39.9\%}（0.121635 $\rightarrow$ 0.073105）。
  \item 相比 TCDR(16,96)，TCDR(32,64) 的 Val $R^2$ 进一步提升 \textbf{+0.0169}，Val Loss 相对下降约 \textbf{18.09\%}。
  \item 两个 TCDR 分支的 Train $R^2$ 接近，但验证指标存在稳定差异，说明提升不仅是拟合增强，更来自跨域迁移稳定性的改善。
\end{enumerate}

\section{主文 One-factor 消融（固定 T32x64）}
来源：\texttt{paper\_materials/tables/paper\_t32\_onefactor\_summary.csv}

\subsection{消融设计说明}
为保证可解释性，采用 one-factor 设计：
\begin{itemize}
  \item 每次只改一个组件；
  \item 其余配置保持一致；
  \item 使用相同协议与相同 seed 集合（9294--9298）；
  \item 统一使用 \texttt{[Best]} 行统计 Train/Val 指标。
\end{itemize}

\begin{table}[htbp]
\centering
\scriptsize
\setlength{\tabcolsep}{3pt}
\renewcommand{\arraystretch}{1.15}
\begin{tabularx}{\linewidth}{L{3.0cm}L{3.8cm}L{3.8cm}X}
\toprule
消融项 & 仅改变的内容 & 保持不变的内容 & 验证目标 \\
\midrule
ModeBase(enrich\_on) &
\texttt{hgat\_net\_feat\_mode=base}；保留 \texttt{enrich\_parasitic\_net\_feat} &
T32x64 主干、拓扑专家、优化器与训练策略 &
确认当前最稳定的 NET 特征路径 \\

EnrichOff(auto) &
关闭 \texttt{enrich\_parasitic\_net\_feat}；\texttt{hgat\_net\_feat\_mode=auto} &
其余全部不变 &
确认 enrich 开关是否提供额外有效信息 \\

ParasiticReplace &
\texttt{hgat\_net\_feat\_mode=}\newline
\texttt{parasitic\_replace} &
其余全部不变 &
验证“替换式寄生特征”与主配置差异 \\

DualMean &
打开 \texttt{hgat\_dual\_readout}；\texttt{hgat\_dual\_merge=mean} &
其余全部不变 &
检验双读出是否带来额外增益 \\

ModeBase+TopoOff &
在 ModeBase 条件下关闭 \texttt{use\_topology\_expert} &
其余全部不变 &
直接量化拓扑残差专家的贡献 \\
\bottomrule
\end{tabularx}
\renewcommand{\arraystretch}{1}
\caption{T32x64 one-factor 设计与控制变量。}
\label{tab:t32_design_cn}
\end{table}

\subsection{T32x64 消融结果}
\begin{table}[htbp]
\centering
\footnotesize
\setlength{\tabcolsep}{4pt}
\begin{tabular}{L{4.3cm}c c c c}
\toprule
方法 & NSeed & Train $R^2$ (mean$\pm$std) & Val $R^2$ (mean$\pm$std) & Val Loss (mean$\pm$std) \\
\midrule
T32x64 + BasePath(mainline) & 5 & 0.9682$\pm$0.0087 & \textbf{0.9298$\pm$0.0073} & \textbf{0.066996$\pm$0.006930} \\
T32x64 + ParasiticReplace & 5 & 0.9694$\pm$0.0080 & 0.9294$\pm$0.0084 & 0.067451$\pm$0.008056 \\
T32x64 + DualMean & 5 & 0.9582$\pm$0.0216 & 0.9234$\pm$0.0175 & 0.073105$\pm$0.016755 \\
T32x64 + ModeBase+TopoOff & 5 & 0.9006$\pm$0.0068 & 0.8720$\pm$0.0034 & 0.122180$\pm$0.003232 \\
\bottomrule
\end{tabular}
\caption{主文 one-factor 消融（固定 T32x64）。}
\label{tab:t32_ablation_cn}
\end{table}

\subsection{T32 one-factor 结果解读}
\begin{enumerate}
  \item \textbf{拓扑专家是核心增益项}：ModeBase(mainline) 与 ModeBase+TopoOff 的 Val $R^2$ 差值为 \textbf{0.0578}（0.9298 vs 0.8720），Val Loss 增量为 \textbf{0.055184}。
  \item \textbf{ParasiticReplace 收益不稳定}：相对 mainline 的 Val $R^2$ 差值仅 \textbf{-0.0005}，说明“替换式”寄生注入未形成稳定优势。
  \item \textbf{DualMean 在当前协议下非收益项}：相对 mainline 的 Val $R^2$ 差值为 \textbf{-0.0064}，并伴随更高的 Val Loss。
  \item \textbf{等价配置需避免重复计功}：EnrichOff(auto) 与 ModeBase(mainline) 在当前实现中落入同一 NET base 路径，指标一致应视为一致性验证，不应在主表中作为两条独立贡献。
\end{enumerate}

\noindent\textbf{组会汇报口径：}
主线方法定义为 \textbf{T32x64 + BasePath(mainline)}。
\texttt{EnrichOff(auto)} 不单列在主表中，仅作为“等价配置一致性验证”：
在当前实现中它会落到同一条 NET base 特征路径，指标与 BasePath(mainline) 一致。

\subsection{为什么主线有优势}
\begin{enumerate}
  \item 关闭拓扑专家后，Val $R^2$ 从 0.9298 降到 0.8720（下降约 0.058），说明拓扑条件残差是核心增益来源。
  \item BasePath 与 EnrichOff(auto) 数值一致，说明主线结果稳定，不依赖脆弱开关。
  \item DualMean 未带来稳定增益，建议作为补充分析而非主卖点。
\end{enumerate}

\section{补充消融（T16x96，对照）}
来源：\texttt{paper\_materials/tables/paper\_notest\_component\_summary.csv}

\begin{table}[htbp]
\centering
\footnotesize
\setlength{\tabcolsep}{4pt}
\begin{tabular}{L{4.0cm}c c c c}
\toprule
方法 & NSeed & Train $R^2$ (mean$\pm$std) & Val $R^2$ (mean$\pm$std) & Val Loss (mean$\pm$std) \\
\midrule
T16x96 + ModeBase(enrich\_on) & 5 & 0.9548$\pm$0.0122 & \textbf{0.9160$\pm$0.0139} & \textbf{0.080192$\pm$0.013224} \\
T16x96 + EnrichOff(auto) & 5 & 0.9548$\pm$0.0122 & \textbf{0.9160$\pm$0.0139} & \textbf{0.080192$\pm$0.013224} \\
T16x96 + ParasiticReplace & 5 & 0.9562$\pm$0.0118 & 0.9151$\pm$0.0113 & 0.081077$\pm$0.010834 \\
T16x96 + DualMean & 5 & 0.9586$\pm$0.0085 & 0.9065$\pm$0.0168 & 0.089249$\pm$0.016086 \\
T16x96 + ModeBase+TopoOff & 5 & 0.9006$\pm$0.0068 & 0.8720$\pm$0.0034 & 0.122180$\pm$0.003232 \\
\bottomrule
\end{tabular}
\caption{补充消融（固定 T16x96，仅作对照）。}
\label{tab:t16_ablation_cn}
\end{table}

\section{实验局限与口径边界}
\begin{itemize}
  \item 当前协议为 \textbf{train:val=1:5 且无独立测试集}，因此本节结论用于模型选择与机制验证，不用于最终泛化主张。
  \item Val 指标最优不等价于跨分布最终最优；后续需在恢复 test 或 cross-group CV 后再给出最终泛化结论。
  \item 当前表格未加入显著性检验（如 paired seed t-test / Wilcoxon），主文投稿前建议补齐。
\end{itemize}

\section{当前可用于组会与写作的结论}
\begin{enumerate}
  \item 在 table\_group + 1:5 + 无测试集协议下，TCDR(32,64) 在 Val 指标最优。
  \item 主线应写为 T32x64 + BasePath(mainline)；EnrichOff(auto) 仅作为等价验证。
  \item Topology expert 是关键贡献组件，关闭后性能明显退化。
  \item 当前结论是阶段性结论；后续恢复独立泛化评估（test 或交叉验证）。
\end{enumerate}

\section{复现入口}
\subsection{T32 主文消融批量运行}
\begin{center}
\texttt{powershell -ExecutionPolicy Bypass -File}\\
\path{ICCAD2026_Changxin/paper_materials/tables/run_step5a_t32_onefactor.ps1}
\end{center}

\subsection{T32 消融结果统计}
\begin{center}
\texttt{powershell -ExecutionPolicy Bypass -File}\\
\path{ICCAD2026_Changxin/paper_materials/tables/build_step5a_t32_onefactor_tables.ps1}
\end{center}

\end{document}
